% Generated draft: requires mathematical and source review.
\documentclass[11pt]{article}
\usepackage[margin=25mm]{geometry}
\usepackage{fontspec}
\setmainfont{Latin Modern Roman}
\usepackage{amsmath,amssymb,mathtools,unicode-math}
\setmathfont{Latin Modern Math}
\usepackage{enumitem,needspace,xurl,hyperref}
\setlength{\parindent}{0pt}
\setlength{\parskip}{5pt}
\setlength{\emergencystretch}{3em}
\newcommand{\N}{\mathbb{N}}
\newcommand{\Z}{\mathbb{Z}}
\newcommand{\Q}{\mathbb{Q}}
\newcommand{\R}{\mathbb{R}}
\newcommand{\C}{\mathbb{C}}
\newcommand{\F}{\mathbb{F}}
\newcommand{\A}{\mathbb{A}}
\newcommand{\PP}{\mathbb{P}}
\newcommand{\E}{\mathbb{E}}
\newcommand{\eps}{\varepsilon}
\newcommand{\dd}{\,\mathrm{d}}
\begin{document}
\section*{UnsolvedMath ID: 5000003}
\textbf{Source problem code:} catalogue entry
\subsection{Definitions and mathematical statement}
Let $n\geq3$ be an integer and fix a regular planar $n$-gon. Unfold a billiard trajectory by reflecting the current polygon across each side that the trajectory encounters, so that the reflected trajectory continues as a straight segment. A short trajectory joins two vertices and meets no vertex between its endpoints. A reachable point is an endpoint vertex arising in such a short-trajectory unfolding, in the source's sense.

For an oriented short trajectory, retain the endpoint angles $\alpha,\beta$ of Definition 2.1 in the source. The record says these give $n-2$ trajectory types $A_k$, with $k$ taken modulo $n-2$, and that
\[
A_0\quad\text{is characterized by}\quad\beta-\alpha=\frac{2\pi}{n}.
\]
The angle conventions and the other types are not reproduced in the record.

Writing $\mathcal R_n$ for the source's reachable points and $\mathcal G_n$ for its reachable $n$-gons, Conjecture 2.4 asserts
\[
\forall q\in\mathcal R_n,\qquad
\#\{Q\in\mathcal G_n:q\text{ is a vertex of }Q\}=\infty.
\]
Equivalently, for every positive integer $m$ there should be $m$ pairwise distinct reachable polygons having $q$ as a vertex. This is stronger than merely producing arbitrarily many reflection words that describe the same polygon.

The source record does not define ``reachable $n$-gon'', fully specify the initial unfolding conventions, or reproduce Definition 2.1. One cannot assume that every reflected polygon, or every polygon whose vertices are reachable, has the required reachability property. The displayed quantified assertion preserves the exact target using the source's classes; their missing definitions are needed for a self-contained statement.
\subsection{Short English statement}
Every endpoint reachable by a vertex-to-vertex billiard unfolding in a regular $n$-gon should be a vertex of infinitely many distinct reachable $n$-gons. The source's polygon-reachability and endpoint-angle definitions need to be supplied.
\subsection{Sources}
\begin{itemize}
\item Local UnsolvedMath registry, permanent ID 5000003, exact source code catalogue entry, ``Reachable points lie on reachable polygons''; proposed by Dmitry Fuchs, 2020. Registry status open, classification OPEN-TRIAGE.
\item Dmitry Fuchs, \emph{Billiard Trajectories in Regular Polygons and Foliations by Closed Geodesics on Surfaces} (2020), Conjecture 2.4 and Definition 2.1, \url{https://amj.math.stonybrook.edu/pdf-Springer-final/020-0170.pdf}. Supplied pdf-text extraction accessed July 29, 2026. No PDF was accessed for this draft.
\item The supplied summary places the conjecture in the family 2.3--2.7, reports extensive computational evidence, and says no later resolution was located in its web search through 2026. Its suggestion that resolution would require a full classification is a supplied research assessment, not a proved necessity or an extra task.
\item Provenance flags status and rights NEEDS\_REVIEW: the paper presents the item as open, no later authoritative resolution check is recorded, and public access did not establish redistribution permission.
\end{itemize}
\end{document}
